\ifdefined\seminarhandout
\documentclass[aspectratio=169,10pt,handout]{beamer}
\else
\documentclass[aspectratio=169,10pt]{beamer}
\fi
\usepackage{nesbeamer}
\title{Seminar 1: Preferences and utility}
\author{Arseniy Stolyarov}
\begin{document}
\begin{frame}{Seminar 1 · Preferences and utility}
\key{Microeconomics}\par\bigskip
Arseniy Stolyarov\\PhD in Economics, University of Wisconsin--Madison\par\bigskip
Relations, contours, representations and counterexamples.
\end{frame}
\begin{frame}{The route}
\begin{enumerate}
\item Classify relations before writing utility.
\item Defend a true claim; break a false one.
\item Separate utility curvature from preference geometry.
\item Distinguish one purchase from the full choice set.
\item Put assumptions back into theorems.
\end{enumerate}
\end{frame}
\begin{frame}{What counts as an answer?}
\question{``False, because the preferences are strange'' is not a counterexample.}
\uncover<2->{\begin{itemize}
\item State the domain and the relation or function.
\item Name the bundles, sequence or menu.
\item Check the hypothesis.
\item Exhibit the precise failed inequality or implication.
\end{itemize}}

\end{frame}
\begin{frame}{1A · Coordinatewise comparison}
On $X=\R_+^2$,
\[x\pref_A y\iff x_1\ge y_1\ \text{and}\ x_2\ge y_2.\]
\question{Complete? Transitive? Continuous?}
\uncover<2->{Incomplete: $(1,0)$ and $(0,1)$.\par
Transitive: chain the two coordinate inequalities.\par
Continuous: two weak inequalities survive limits.}

\end{frame}
\begin{frame}{1B · Replace AND by OR}
\[x\pref_B y\iff x_1\ge y_1\ \text{or}\ x_2\ge y_2.\]
\question{Complete? Transitive? Continuous?}
\uncover<2->{Complete: failing both inequalities makes the reverse comparison true.\par
Intransitive: $x=(1,0)$, $y=(0,2)$, $z=(2,1)$ satisfy $x\pref_B y\pref_B z$ but $x\not\pref_B z$.\par
Continuous: its graph is a finite union of closed sets.}

\end{frame}
\begin{frame}{1C · A tiny tie-breaker?}
Compare $x_1+x_2$ first; when totals tie, compare $x_1$.
\question{Can $x_1+x_2+\varepsilon x_1$, for some fixed $\varepsilon>0$, represent this ranking?}
\uncover<2->{No. Take $a=(0,1+\delta)$, $b=(1,0)$ and $0<\delta<\varepsilon$.\par
The primary total ranks $a$ higher; the weighted sum ranks $b$ higher.\par
The relation is complete and transitive, but discontinuous.}

\end{frame}
\begin{frame}{1C · Find the broken limit}
\question{Find $x^n\pref_C y^n$ whose limits reverse the comparison.}
\uncover<2->{\[
x^n=(0,1+1/n),\qquad y^n=(1,0).
\]
Every $x^n$ has the larger total. At the limit, totals tie and $(1,0)$ wins the tie-breaker.}

\end{frame}
\begin{frame}{2 · True or false?}
\begin{enumerate}
\item Incomparability means indifference.
\item Utility representation implies rationality.
\item Discontinuous utility implies discontinuous preferences.
\end{enumerate}
\uncover<2->{\begin{enumerate}
\item False: neither comparison versus both comparisons.
\item True: inherit completeness and transitivity from $\ge$.
\item False: $x_1+\mathbf1\{x_1\ge1\}$ represents the same continuous ranking as $x_1$.
\end{enumerate}}

\end{frame}
\begin{frame}{2 · Three more claims}
\begin{enumerate}
\item Convex utility implies convex preferences.
\item Increasing transformations preserve convex preferences, but need not preserve concavity of utility.
\item LNS implies utility is unbounded above.
\end{enumerate}
\uncover<2->{\begin{enumerate}
\item False: $u=x_1^2+x_2^2$; mix $(2,0)$ and $(0,2)$.
\item True: same rankings and upper contour sets.
\item False: $u=1-e^{-(x_1+x_2)}$ is bounded and LNS.
\end{enumerate}}

\end{frame}
\begin{frame}{3 · A kink is not a pathology}
\[u(x)=\min\{2x_1,x_2\},\qquad X=\R_+^2.\]
\question{Draw the indifference set and the upper contour set through $(1,2)$.}
\uncover<2->{\begin{columns}
\column{.52\textwidth}
\begin{tikzpicture}\begin{axis}[nes,width=6.1cm,height=3.8cm,xmax=4,ymax=5,xtick={0,1,2,3,4},ytick={0,2,4}]
\addplot[draw=none,fill=NESlight,fill opacity=.35] coordinates {(1,2)(4,2)(4,5)(1,5)}\closedcycle;
\addplot[NESblue] coordinates {(1,5)(1,2)(4,2)};
\addplot[only marks,NESred] coordinates {(1,2)};
\end{axis}\end{tikzpicture}
\column{.44\textwidth}
$u=2$: two rays.\par
$u\ge2$: $x_1\ge1$ and $x_2\ge2$.
\end{columns}}

\end{frame}
\begin{frame}{3 · Classify the preferences}
\question{Rational? Continuous? Convex? Strictly convex? Monotone? Strongly monotone? LNS?}
\uncover<2->{Rational, continuous and convex.\par
Not strictly convex: $(1,3)\sim(1,4)$ and their midpoint.\par
Monotone and LNS: increase both goods slightly.\par
Not strongly monotone: $(1,2)\sim(2,2)$.}

\end{frame}
\begin{frame}{3 · Transform the index}
\[v=e^u,\qquad w=\mathbf1\{u\ge2\}.\]
\question{Same preferences? Concave utility?}
\uncover<2->{$u$ is concave: the minimum of these two linear functions is concave.\par
$v$ preserves preferences but is not concave: along $(t,2t)$ it equals $e^{2t}$.\par
$w$ loses strict comparisons, e.g. between $(1,2)$ and $(2,4)$. It is not another representation.}

\end{frame}
\begin{frame}{4 · A choice changes when the menu expands}
Selected $a$ from $\{a,b\}$; selected $b$ from $\{a,b,c\}$.
\question{Does this refute a fixed rational ranking?}
\uncover<2->{Not if each record is only one selected maximiser.\par
$a\sim b\succ c$ permits both selections.\par
We need to know whether unselected alternatives were also best.}

\end{frame}
\begin{frame}{4 · Now these are full best-alternative sets}
\[C(\{a,b\})=\{a\},\qquad C(\{a,b,c\})=\{b\}.\]
\question{Can a rational ranking generate these correspondences?}
\uncover<2->{No. The pair implies $a\succ b$, whereas the triple implies $b\succ a$.\par
The interpretation of the data changes the conclusion.}

\end{frame}
\begin{frame}{5 · An optimum cannot leave slack under LNS}
\[B=\{x\ge0:p\cdot x\le m\},\qquad p\gg0,\ m>0.\]
Assume a best bundle $x^*$ exists.
\question{Prove budget exhaustion without differentiating.}
\uncover<2->{If $s=m-p\cdot x^*>0$, choose $\delta=s/(2\|p\|)$.\par
Every consumption bundle within distance $\delta$ remains affordable.\par
LNS gives a strictly better nearby bundle: contradiction.}

\end{frame}
\begin{frame}{5 · Does the converse hold?}
\[u=-[(x_1-2)^2+(x_2-2)^2],\quad p=(1,1),\quad m=1.\]
\question{Every optimum exhausts this budget. Does that prove LNS?}
\uncover<2->{No. Below this budget line, slightly increasing either good improves utility.\par
But LNS fails at $(2,2)$, the global bliss point.\par
A statement about one budget cannot establish a global property.}

\end{frame}
\begin{frame}{5 · A failed counterexample is useful}
On $x_1+x_2\le1$, lexicographically maximise $x_1$, then minimise $x_2$.
\question{Does this show that LNS need not give existence?}
\uncover<2->{No: $(1,0)$ is a maximiser.\par
A valid repair uses $X=[0,1]$, $r(x)=x$ for $x<1$, $r(1)=-1$.\par
It is LNS but has no maximum, despite compactness. Continuity fails.}

\end{frame}
\begin{frame}{6 · Reserve: discretise lexicographic preferences}
\[X_N=\{0,1/N,\ldots,1\}^2,\qquad U_N=Mx_1+x_2.\]
\question{Which $M$ works? Why does the same trick fail on the square?}
\uncover<2->{$M>N$: the smallest primary gain outweighs the largest secondary loss.\par
On $[0,1]^2$, primary gains can be arbitrarily small.\par
In fact no real-valued representation exists there: use disjoint utility intervals and countability of the rationals.}

\end{frame}
\begin{frame}{What to carry into the homework}
\begin{itemize}
\item Quantifiers and domains are part of every claim.
\item A representation is ordinal; a curvature property need not be.
\item Convex preferences mean convex \emph{upper} contour sets.
\item Keep full best sets separate from selected purchases.
\item Existence, uniqueness and budget exhaustion need different assumptions.
\end{itemize}
\end{frame}
\end{document}
